\section{A distribuição inicial e sua evolução}
\label{sec:distribuicao}

\subsection{O instante do surto (\texorpdfstring{$t=0^+$}{t=0+}): divisor capacitivo}

No primeiro instante após a chegada da frente, os indutores bloqueiam a corrente e o enrolamento é uma
rede puramente capacitiva. Resolvendo essa rede com o neutro aterrado, a tensão se distribui segundo a
solução clássica \cite{greenwood1991}:

\begin{equation}
\frac{v(x)}{V_0} = \frac{\sinh\!\big(\adist(1-x)\big)}{\sinh(\adist)},
\qquad x \in [0,1],
\label{eq:sinh}
\end{equation}

em que \(x=0\) é a entrada (energizada) e \(x=1\) é o neutro aterrado. A Figura~\ref{fig:dist-inicial}
mostra esse perfil, calculado pelo \emph{modelo} (não desenhado à mão), para vários valores de
\(\adist\): quanto maior \(\adist\), mais a tensão se acumula nas primeiras seções e mais a curva se
afasta da reta uniforme.

\staticfigure{0.92}{distribuicao_inicial.png}{Distribuição inicial de tensão \(v(x)/V_0\) em
\(t=0^+\), com neutro aterrado, obtida do modelo distribuído (\codefile{initial\_voltage\_distribution})
para \(\adist = 1, 3, 5, 8\). A reta tracejada é a distribuição uniforme ingênua (\(\adist\to 0\)).
Para a máquina estudada, \(\adist=5\).}{fig:dist-inicial}

A não uniformidade é medida pelo \emph{fator de concentração na entrada}, o gradiente local na primeira
espira em relação ao gradiente médio:

\begin{equation}
\left.\frac{\mathrm{d}}{\mathrm{d}x}\frac{v}{V_0}\right|_{x=0}
   = \adist\coth(\adist) \approx \adist = 5 .
\label{eq:concentracao}
\end{equation}

Ou seja: para \(\adist=5\), a primeira fração do enrolamento é solicitada por cerca de \textbf{cinco
vezes} a tensão que veria se a distribuição fosse uniforme. É esse o mecanismo que estressa o isolamento
de entrada de máquinas e transformadores.

\subsection{A evolução no tempo: do perfil concentrado ao uniforme}

A distribuição da Equação~\eqref{eq:sinh} é apenas o ponto de partida. À medida que a corrente começa a
circular pelos indutores, a energia se redistribui e o perfil \emph{relaxa} de uma forma concentrada
para uma distribuição quase uniforme, presa a zero no neutro aterrado. Durante essa transição surgem
oscilações internas e sobretensões: a simulação no tempo indica pico em torno de \SI{1108}{V} no nó a
25\,\% do enrolamento (\(N_5\)) — acima do \SI{1}{kV} aplicado na entrada. A Seção~\ref{sec:animacoes}
mostra essa evolução em animação.

A leitura de engenharia é que o \emph{envelope} de tensão ao longo do enrolamento — e não o regime
permanente — dimensiona o isolamento: as primeiras espiras concentram a solicitação no instante do
surto, e nós internos podem exceder a tensão de entrada durante a redistribuição.
